% Generated draft: requires mathematical and source review.
\documentclass[11pt]{article}
\usepackage[margin=25mm]{geometry}
\usepackage{fontspec}
\setmainfont{Latin Modern Roman}
\usepackage{amsmath,amssymb,mathtools,unicode-math}
\setmathfont{Latin Modern Math}
\usepackage{enumitem,needspace,xurl,hyperref}
\setlength{\parindent}{0pt}
\setlength{\parskip}{5pt}
\setlength{\emergencystretch}{3em}
\newcommand{\N}{\mathbb{N}}
\newcommand{\Z}{\mathbb{Z}}
\newcommand{\Q}{\mathbb{Q}}
\newcommand{\R}{\mathbb{R}}
\newcommand{\C}{\mathbb{C}}
\newcommand{\F}{\mathbb{F}}
\newcommand{\A}{\mathbb{A}}
\newcommand{\PP}{\mathbb{P}}
\newcommand{\E}{\mathbb{E}}
\newcommand{\eps}{\varepsilon}
\newcommand{\dd}{\,\mathrm{d}}
\begin{document}
\section*{UnsolvedMath ID: 1200027}
\textbf{Source problem code:} catalogue entry
\subsection{Definitions and mathematical statement}
Let \(\Gamma\) be an infinite group and \(X=\operatorname{Cay}(\Gamma,S)\) a connected Cayley graph, using the usual undirected locally finite convention: \(S\) is a finite symmetric generating set with \(1\notin S\), and edges join \(g\) to \(gs\). A perfect matching is a set \(M\subseteq E(X)\) such that each vertex is incident to exactly one edge of \(M\). Write \(\mathcal M(X)\subseteq\{0,1\}^{E(X)}\) for the space of perfect matchings, with its product Borel structure. Left translations of \(\Gamma\) act on \(\mathcal M(X)\).

An invariant random perfect matching is a Borel probability measure \(\mu\) on \(\mathcal M(X)\) satisfying
\[
\mu(\gamma B)=\mu(B)\qquad
(\gamma\in\Gamma,\ B\subseteq\mathcal M(X)\text{ Borel}).
\]
The question asks whether every infinite connected Cayley graph admits such a measure. Invariance here is under the defining group action, as explicitly specified in the supplied research summary.

That summary records an affirmative answer and adds that nonamenable Cayley graphs admit such a matching as a factor of independent identically distributed labels: a measurable \(\Gamma\)-equivariant function of independent vertex labels, producing a perfect matching almost surely. This is a supplied literature claim, not a proof or a newly verified theorem.
\subsection{Short English statement}
Can every infinite connected Cayley graph be perfectly matched at random so that its distribution is unchanged by group translations? The supplied research summary answers yes and reports a factor-of-i.i.d. construction in the nonamenable case.
\subsection{Sources}
\begin{itemize}
\item Local UnsolvedMath registry, permanent ID \(1200027\), exact source code \texttt{catalogue entry}, ``Some Questions --- Question 27.''
\item Abert, \emph{Some questions} (2010), Question 27; \url{https://www.renyi.hu/~abert/questions.pdf}; manual transcription; recorded access 2026-07-29.
\item Supplied research summary: Cs\'oka--Lippner--Pikhurko, \emph{Groups, Geometry, and Dynamics} \textbf{11} (2017), 211--243; arXiv:1211.2374. The supplied citation has not been independently checked here.
\end{itemize}
\end{document}
